\section{Conclusion}


This thesis marks one of the first applications of DEQNs to an IAM. It details the methodology required to adapt ACE to the DEQN framework, offering a more comprehensive treatment than previous studies. Given ACE's general structure, the methods and derivations in this thesis should be easily extended to other IAMs, facilitating broader adoption in future research.

The thesis conducts two experiments utilizing the mapping of ACE to the DEQN framework, with the DICE 2016 and DICE 2023 production functions. With the corrected equilibrium conditions, the policy network recovers the analytical solution closely in both experiments. The estimated consumption rate stays within 0.3 percentage points of the constant analytical rate, the estimated abatement rate follows the analytical path to full abatement with a mean deviation of 0.1 to 0.2 percentage points, and the costate variables, including the carbon costate that is four orders of magnitude smaller than the capital costate, are recovered to within 1.5\%. As a result, the SCC implied by the policy network and the simulated paths of temperature and atmospheric carbon are close to their analytical counterparts.

The revision of this thesis also carries a methodological lesson. In the submitted version, the envelope condition for capital omitted the direct effect of capital on utility, and the first-order conditions counted the effect of today's controls on the future value twice. The policy network then converged towards a wrong fixed point, and I attributed the poor results to the normalization of inputs. Because ACE has an analytical solution, the conditions can be verified directly by evaluating them at the known optimum, which reveals such errors before any network is trained. For models without a closed-form solution, checking the equilibrium conditions on a simplified version with a known solution is therefore a valuable step. A second lesson is that the conditions should be weighted so that each carries comparable weight in the policy objective: without scaling, the conditions on the carbon cycle are several orders of magnitude smaller than the others and barely identify the SCC.

This work establishes a connection with RL by comprehensively discussing concepts common to both DEQNs and the broader RL literature. Contrary to previous research that classified this algorithm as unsupervised learning, this thesis argues that the DEQN algorithm should be categorized as model-based RL. This reclassification provides a foundation for the potential integration of more sophisticated RL algorithms in future economic research. With a correctly specified objective, natural next steps are to relax the assumptions that make ACE analytically tractable, such as the treatment of business-as-usual emissions as given and the logarithmic utility, and to introduce uncertainty, where the advantages of DEQNs over traditional solution methods become relevant.
